\chapter{Riwayat dan Masalah Umum}

Setiap pekerjaan di Pandora tersimpan sebagai folder hasil di riwayat akun Anda. Riwayat memudahkan Anda membuka kembali hasil lama, melanjutkan pekerjaan yang belum selesai, dan menghapus hasil yang tidak diperlukan. Bab ini membahas cara mengelola riwayat dan cara menangani masalah yang paling sering muncul. Masalah yang tidak tercantum di sini dilaporkan kepada pengelola bersama langkah yang sudah Anda coba.

\begin{indikator}
Setelah mempelajari bab ini, peserta mampu:
\begin{enumerate}
  \item membuka kembali hasil lama dari riwayat dan melanjutkan pekerjaan yang belum selesai;
  \item menghapus hasil dengan aman setelah memahami akibatnya;
  \item mengenali lima masalah umum dan mengambil langkah pertama untuk masing-masing.
\end{enumerate}
\end{indikator}

\section{Membuka dan Melanjutkan Riwayat}
Panel riwayat menampilkan semua folder hasil milik Anda, dengan nama yang mengikuti pola \texttt{01\_nama}, \texttt{02\_nama}, dan seterusnya. Klik sebuah folder untuk membuka hasilnya kembali, termasuk transkrip, rapi AI, notulensi, dan dokumen bila sudah dibuat.

Bila sebuah folder belum memiliki notulensi atau dokumen, tombol untuk membuatnya tetap tersedia setelah folder dibuka. Dengan cara ini, Anda tidak perlu mentranskripsi ulang rekaman yang sama.

\section{Menghapus Riwayat}
Menghapus folder hasil berarti menghapus semua berkas di dalamnya, yaitu transkrip, hasil rapi, notulensi, dan dokumen. Pandora meminta konfirmasi sebelum menghapus. Pesan konfirmasi menyebut nama folder yang akan dihapus. Penghapusan tidak dapat dibatalkan.

Sebelum menghapus, pastikan Anda tidak lagi memerlukan hasilnya, dan bila perlu unduh dokumen yang masih dibutuhkan. Jangan menghapus hasil yang masih Anda perlukan sebagai rujukan, misalnya transkrip asli yang menjadi sumber notulensi atau dokumen yang sudah diajukan.

\section{Masalah Umum dan Langkah Pertama}
Tabel~\ref{tab:masalah} merangkum masalah yang paling sering muncul beserta langkah pertama yang bisa Anda lakukan.

\begin{table}[H]
\caption{Masalah umum dan langkah pertama}
\label{tab:masalah}
{\small
\begin{tabularx}{\linewidth}{@{}P{4.2cm}YY@{}}
\toprule
\rowcolor{ugmpanel}\thd{Gejala} & \thd{Kemungkinan penyebab} & \thd{Langkah pertama} \\
\midrule
Tombol AI tidak dapat diklik & Layanan AI belum siap & Periksa status di bagian atas layar, lalu hubungi pengelola \\
Pesan worker sedang digunakan & Transkripsi lain sedang berjalan & Tunggu sampai selesai, lalu coba lagi \\
Berkas tidak dapat diunggah & Format berkas tidak didukung & Ubah ke format yang didukung, lalu unggah ulang \\
Log berulang atau kosong di bagian tertentu & Bagian hening atau suara tidak jelas & Periksa rekaman di menit tersebut sebelum menyimpulkan \\
Notulensi atau dokumen tidak selesai & Layanan AI sibuk atau terhenti & Tunggu, lalu coba lagi. Bila tetap gagal, catat pesan dan hubungi pengelola \\
Naskah ditolak pemeriksa & Cakupan, bentuk, atau keterlacakan tidak terpenuhi & Baca alasan yang ditampilkan, lalu susun ulang sesuai bagian \\
\bottomrule
\end{tabularx}
}
\end{table}

Bila masalah masih berlanjut setelah langkah pertama, catat tiga hal: waktu kejadian, pesan yang muncul, dan nama folder hasil yang terkait. Informasi ini membantu pengelola menelusuri masalah tanpa perlu menanyakan ulang.

\begin{ilustrasi}
Pada hari kerja yang padat, Rina mencoba memulai transkripsi rapat kedua, tetapi layar menampilkan pesan bahwa worker sedang digunakan. Ia tidak mengulang berkali-kali. Ia menunggu, lalu mencoba lagi sekitar dua puluh menit kemudian. Setelah itu, ia membuka riwayat dan menghapus folder percobaan pertamanya yang tidak lagi diperlukan, dengan membaca pesan konfirmasi terlebih dahulu.
\end{ilustrasi}

\begin{kegiatan}{Latihan 6. Menelusuri riwayat dan masalah (10 menit)}
Buka riwayat dan pilih satu folder hasil dari latihan sebelumnya. Tentukan apakah notulensi dan dokumen sudah ada di dalamnya. Lalu pilih satu gejala dari Tabel~\ref{tab:masalah} dan tulis langkah pertama yang akan Anda ambil, tanpa benar-benar menghapus apa pun. Keluaran yang diharapkan: daftar kondisi folder hasil dan satu langkah pertama yang Anda pilih.
\end{kegiatan}

\begin{rangkuman}
Riwayat menyimpan semua folder hasil milik akun Anda dan memungkinkan Anda melanjutkan pekerjaan lama. Menghapus folder berarti menghapus seluruh isinya, sehingga Pandora meminta konfirmasi terlebih dahulu. Sebagian besar masalah umum dapat diatasi dengan menunggu, memeriksa status, atau menyusun ulang naskah sesuai alasannya. Rina menunggu giliran transkripsi dan menghapus hasil yang tidak lagi diperlukan. Bab berikutnya menutup modul dengan kembali ke pertanyaan pemantik.
\end{rangkuman}
